\begin{proof}
Notation as in the statement: $D_{\mu} = D_{e_{\mu}}$, $g^{\mu} = \hb^{T}(e^{\mu})$, $g_{\mu} = \hb^{-1}(e_{\mu})$, so that $g^{\mu}\cdot g_{\nu} = \delta^{\mu}_{\nu}$, $D = g^{\mu}D_{\mu}$, and $\partial_{a} = g^{\mu}\,(g_{\mu}\cdot\partial_{a})$ for a physical argument $a$. The operator $D_{\mu}$ is a derivation of the geometric product, since $\Omega(e_{\mu})\times$ is one, and it preserves grades.

\emph{(i).} For a vector $c$, a vector $x$, a vector $y$ and a bivector $B$, the rule $c\cdot(A_{r}\wedge Z) = (c\cdot A_{r})\wedge Z + (-1)^{r}A_{r}\wedge(c\cdot Z)$ gives
\[
x\wedge y\wedge(c\cdot B) = c\cdot(x\wedge y\wedge B) - (c\cdot x)\,y\wedge B + (c\cdot y)\,x\wedge B .
\]
Put $x = \partial_{b}$, $y = \partial_{a}$, $B = M(a\wedge b)$ and sum. The first term is $c\cdot\bigl(\partial_{b}\wedge\partial_{a}\wedge M(a\wedge b)\bigr) = c\cdot\bigl(\partial_{b}\wedge P_{M}(b)\bigr) = c\cdot E_{M}$. The second is $-(c\cdot\partial_{b})\,\partial_{a}\wedge M(a\wedge b) = -\partial_{a}\wedge M(a\wedge c) = -P_{M}(c)$. The third is $(c\cdot\partial_{a})\,\partial_{b}\wedge M(a\wedge b) = \partial_{b}\wedge M(c\wedge b) = -P_{M}(c)$. So $\partial_{b}\wedge\partial_{a}\wedge(c\cdot M(a\wedge b)) = c\cdot E_{M} - 2P_{M}(c)$, and since $M(a\wedge b)\cdot c = -c\cdot M(a\wedge b)$,
\[
C_{M}(c) = 2P_{M}(c) - c\cdot E_{M} .
\]
For a multivector $X_{r}$ of grade $r$, $\partial_{c}\wedge(c\cdot X_{r}) = r\,X_{r}$; $E_{M}$ has grade $4$, so $\partial_{c}\wedge C_{M}(c) = 2E_{M} - 4E_{M} = -2E_{M}$. Substituting $E_{M} = -\tfrac12\partial_{e}\wedge C_{M}(e)$ gives the inversion formula. The grade count $4$ does not depend on the dimension or the signature. The check verifies the three identities on all basis maps and frame vectors in two signatures of dimension four, which is complete there by linearity, and on a sample in dimension five.

For the vector form: pairing the trivector $C_{M}(c)$ with $z\wedge y\wedge x$ antisymmetrises $(M(a\wedge b)\cdot c)\cdot z'$ over the three vectors placed in the slots $a$, $b$, $z'$; since $M(a\wedge b)$ is antisymmetric in $a, b$ and $(B\cdot c)\cdot z = -(B\cdot z)\cdot c$ for a bivector $B$, the result is a nonzero multiple of $c\cdot\bigl(M(x\wedge y)\cdot z + M(y\wedge z)\cdot x + M(z\wedge x)\cdot y\bigr)$. So $C_{M}(c) = 0$ for every $c$ exactly when the cyclic vector sum vanishes for all $x, y, z$. In $\Cl(1,3)$ the check confirms this by exact rank: the cyclic vector identity and the protraction identity have the same solution space.

\emph{(ii).} By the Leibniz rule,
\[
D\wedge\Theta^{\mu} = g^{\lambda}\wedge D_{\lambda}\bigl(g^{\nu}\wedge D_{\nu}g^{\mu}\bigr) = (g^{\lambda}\wedge D_{\lambda}g^{\nu})\wedge D_{\nu}g^{\mu} + g^{\lambda}\wedge g^{\nu}\wedge D_{\lambda}D_{\nu}g^{\mu} = \Theta^{\nu}\wedge D_{\nu}g^{\mu} - g^{\nu}\wedge g^{\lambda}\wedge D_{\lambda}D_{\nu}g^{\mu} .
\]
The outer product $g^{\nu}\wedge g^{\lambda}$ is antisymmetric, so only the commutator survives: by clause (c) of the definition, $[D_{\lambda}, D_{\nu}]g^{\mu} = R(e_{\lambda}\wedge e_{\nu})\times g^{\mu} = R(e_{\lambda}\wedge e_{\nu})\cdot g^{\mu}$ (the commutator product of a bivector with a vector is their inner product). Since $\hb(g_{\lambda}\wedge g_{\nu}) = e_{\lambda}\wedge e_{\nu}$, $R(e_{\lambda}\wedge e_{\nu}) = -\mathcal{R}(g_{\lambda}\wedge g_{\nu})$. Hence
\[
g^{\nu}\wedge g^{\lambda}\wedge D_{\lambda}D_{\nu}g^{\mu} = -\tfrac12\,g^{\nu}\wedge g^{\lambda}\wedge\bigl(\mathcal{R}(g_{\lambda}\wedge g_{\nu})\cdot g^{\mu}\bigr) = -\tfrac12\,\partial_{b}\wedge\partial_{a}\wedge\bigl(\mathcal{R}(a\wedge b)\cdot g^{\mu}\bigr) = -\tfrac12\,C_{\mathcal{R}}(g^{\mu}),
\]
which is (ii). The second derivatives of $\hb$ that occur in $D_{\lambda}D_{\nu}g^{\mu}$ are symmetric and drop out of the commutator, which is why $\hb$ twice and $\Omega$ once continuously differentiable suffice. The check verifies (ii) exactly on two polynomial fields with torsion.

\emph{(iii).} Vanishing torsion, clause (f) of the definition, reads $D\wedge\hb^{T}(a) = \hb^{T}(\nabla\wedge a)$ for vector fields $a$. For $a = f\,e^{\mu}$ with $f$ a scalar field both sides change by $\hb^{T}(\nabla f)\wedge g^{\mu}$ when $f$ is brought outside ($\hb^{T}$ is an outermorphism and $D f = \hb^{T}(\nabla f)$), so the equation holds for all $a$ exactly when it holds for the constant vectors $e^{\mu}$, where it reads $\Theta^{\mu} = 0$. Then the right side of (ii) vanishes, $C_{\mathcal{R}}(g^{\mu}) = 0$ for every $\mu$, so $C_{\mathcal{R}} = 0$ by linearity ($\{g^{\mu}\}$ is a basis), and (i) gives $P_{\mathcal{R}} = 0$, with the cyclic vector identity as its equivalent form. Finally $\partial_{a}\,\mathcal{R}(a\wedge b) = \partial_{a}\cdot\mathcal{R}(a\wedge b) + \partial_{a}\wedge\mathcal{R}(a\wedge b) = \mathcal{R}(b)$. The check verifies each statement on six torsion-free fields, with $\Omega$ solved exactly from the torsion equation for four polynomial coframes, and that the two controls with torsion of (ii) have $P_{\mathcal{R}} \neq 0$.

\emph{(iv).} The map $M \mapsto P_{M}$ is linear from the maps of bivectors, a space of dimension $n^{2}$ with $n = d(d-1)/2$, to the trivector-valued linear functions of a vector. The check computes its rank exactly: $3$, $16$, $16$, $50$ in signatures $(1,2)$, $(1,3)$, $(2,2)$, $(1,4)$, so the kernel has dimension $6$, $20$, $20$, $50$, which is $d^{2}(d^{2}-1)/12$; adjoining the symmetry conditions $\langle B_{i}M(B_{j})\rangle_{0} = \langle M(B_{i})B_{j}\rangle_{0}$ on the basis bivectors does not raise the rank, so every map of the kernel is symmetric; and the symmetric maps, of dimension $n(n+1)/2$, exceed the kernel by $0$, $1$, $1$, $5$ dimensions. For a symmetric map the Ricci map is symmetric: by $(a\wedge c)\cdot B = a\cdot(c\cdot B)$, $a\cdot M(b) = a\cdot(\partial_{c}\cdot M(c\wedge b)) = \langle (a\wedge\partial_{c})\,M(c\wedge b)\rangle_{0} = \langle M(a\wedge\partial_{c})\,(c\wedge b)\rangle_{0} = c\cdot\bigl(b\cdot M(a\wedge\partial_{c})\bigr)$, and exchanging the roles of $c$ and $\partial_{c}$ in the sum over a frame and its reciprocal, this is $\partial_{c}\cdot(b\cdot M(a\wedge c)) = b\cdot(\partial_{c}\cdot M(c\wedge a)) = b\cdot M(a)$. In $\Cl(1,3)$ the map $B \mapsto I B$ is symmetric because $I$ commutes with bivectors, and since $I$ anticommutes with vectors, $P(c) = \partial_{a}\wedge(I\,a\wedge c) = -I\,\partial_{a}\cdot(a\wedge c) = -3\,I c \neq 0$. With (iii) this gives the symmetry of the curvature operator of a torsion-free field in the dimensions computed; in every dimension it is the classical pair symmetry of an algebraic curvature tensor, cited and not proved here.
\end{proof}
